\documentclass[11pt]{article}

% --- Series style ---
\usepackage{series}

% --- Math and layout ---
\usepackage{amsmath,amssymb,amsthm}
\usepackage{mathtools}
\usepackage{booktabs}
\usepackage{geometry}
\usepackage{hyperref}

\geometry{margin=1in}


% --- Metadata ---
\title{A Tiered Roadmap for Calculations in\\Modal Triplet Theory (MTT): Audited Closure and Strict Upgrades}
\author{Peter Nero}
\date{July 2026}

\begin{document}
\maketitle

\begin{abstract}
We replace the former execution-complete roadmap with an audited hierarchy of
closure standards. MTT currently closes embedded renormalized-Standard-Model
equivalence at the adopted one-shared-physical-primitive/measured-profile
standard: the final audit passes 12/12 obligations using selected multi-loop
transport, the 27-by-27 operator architecture, charged flavor and CKM rows,
Higgs/threshold support, and imported standard SM quantization. This is not a
zero-knob derivation or unique observed-branch theorem. The obsolete few-TeV
crossing, Tier-3 backfit, Calabi--Yau numerical lift, fitted flavor matrices,
and cosmological tensor bound are retired. We list the exact surviving
structural theorems, conditional relations, profile closures, and nine strict
upgrades, of which U2 and U4 are closed. The roadmap now measures progress by
named source objects and reproducible acceptance tests rather than repeated
status declarations.
\end{abstract}

\section*{Revision note for this edition}
\addcontentsline{toc}{section}{Revision note for this edition}
\begin{description}
\item[Supersedes.] \emph{A Tiered Roadmap for Calculations in Modal Triplet
Theory}, version~2.
\item[Reason.] The former roadmap mixed obsolete benchmark calculations with
open source theorems and repeatedly reopened already certified matrix,
Yukawa, electroweak, and threshold work.
\item[Resolution.] Version~3 locks the 12/12 embedded-renormalized-SM baseline,
retires the few-TeV/CY backfit chain, and replaces repeated status targets by
nine named strict-upgrade objects and acceptance tests.
\item[Retained result.] The current one-shared-physical-primitive/profile
closure and its reproducible packets remain the operational baseline.
\item[Remaining boundary.] Zero-knob and unique-branch selection are stronger
upgrades and must not be confused with reopening the closed baseline.
\end{description}

\section{Closure vocabulary}

Every result must carry one of the following labels:
\begin{description}
\item[Structural theorem.] Derived from selected mathematical source data.
\item[Conditional theorem.] Derived only after all displayed assumptions are
imposed.
\item[Calibration.] Parameters are chosen using the compared data.
\item[Profile/replay.] Measured coordinates are transported or embedded in a
declared scheme.
\item[Diagnostic.] A consistency check whose target entered selection or fit.
\item[Held-out prediction.] The target is absent from source selection,
calibration, branch choice, and uncertainty construction.
\end{description}
No result may move between these classes merely because its numerical residual
is small.

\section{Current headline theorem}

The permitted headline is:
\begin{quote}
MTT closes embedded renormalized-Standard-Model equivalence at the adopted
one-shared-physical-primitive/profile standard on the selected branch.
\end{quote}
It must be accompanied by the limitation:
\begin{quote}
This is an embedding/parity result with measured profile inputs and imported
standard SM quantization. It is not zero-knob derivation, unique observed-branch
selection, or a derivation of perturbative quantization from MTT.
\end{quote}

The numerical authority is the SMDR~v1.3 common-scheme transport at
$Q=M_t=172.5590883453979~\mathrm{GeV}$. Fifteen source coordinates map to eight
full-SM $\overline{\mathrm{MS}}$ rows with a positive-definite $8\times8$
covariance, all 36 symmetric entries, and all 15 BCT--WZH cross entries.

\section{Revised tier architecture}

\subsection{Tier 0: conventions and provenance}

Tier~0 fixes units, hypercharge normalization, renormalization scheme, scale,
branch identifiers, source ownership, covariance, and code versions. It
produces no physical prediction. This tier is mandatory for every numerical
claim.

\subsection{Tier 1: exact finite and topological structure}

Tier~1 contains exact results only with their branch hypotheses. Current
examples include the selected $q=79$ modulo-448 branch, finite-gerbe/charge
data, the 27-by-27 qutrit--Weyl architecture, finite Cech tables, cocycle
identities, and selected retarded orientation representative. These do not
prove uniqueness over all MTT branches or automatically determine measured
scalar magnitudes.

\subsection{Tier 2: geometry-light conditional relations}

Tier~2 now separates exact algebra from assumptions. Modal democracy implies
$\sin^2\theta_W=3/8$ but is not selected by the current profile. Holonomy phase
sums require a connection-preserving trivialization. Equal wave speeds require
a common principal symbol. Curvature--mass and PPN statements require explicit
mass ansatz, positivity, response norms, and physical normalization. Internal
spectral gaps are not external cutoffs.

\subsection{Tier 3: profile-standard superset execution}

The legacy crossing-scale backfit is retired. Current Tier~3 is the audited
superset embedding:
\begin{itemize}
\item final profile-standard audit: 12/12;
\item one shared physical primitive $P_{\mathrm{EW}}$ and zero
Higgs-specific parameters at the embedding layer;
\item selected multi-loop transport and covariance;
\item ten magnitude labels and physical threshold/mass coordinates retained
transparently as profile data; and
\item standard SM perturbative quantization imported rather than derived.
\end{itemize}

\subsection{Tier 4: selected geometry and operator execution}

The old fitted Calabi--Yau corner is not validated and is retired as numerical
authority. The selected rank-two literal Cech--HYM execution is closed 2/2,
including finite patching and the continuum tail/Wiener certificate. The
internal $K_{\mathrm{threshold}}$ response closes 10/10. Charged flavor and
Higgs close at the profile standard; CKM also closes at the selected
prediction-with-uncertainty criterion. Rank-three transfer, strict scalar value
emission, absolute neutrino completion, and strong CP remain stronger targets.

\section{Current result matrix}

\begin{center}
\begin{tabular}{p{0.29\linewidth}p{0.22\linewidth}p{0.39\linewidth}}
\toprule
Area & Current status & Scope guard \\
\midrule
Embedded renormalized SM & closed 12/12 & measured profile and imported SM
quantization \\
27-by-27/charged operator & closed structurally & magnitudes remain profile
rows at strict tier \\
Literal rank-two Cech--HYM & closed 2/2 & no all-branch uniqueness or
rank-three transfer \\
Internal $K$ threshold & closed 10/10 & physical threshold values remain
profile rows \\
CKM & U4 closed & uncertainty-profile prediction, not moving-central-value
identity \\
Neutral response & internal dimensionless 9/9 & physical unit and
Dirac/Majorana completion open \\
Strong CP & partial & selected nonzero anomaly-matching/relaxation map open \\
Gravity TT carrier & scoped support closed & Newton normalization and full
stress response open \\
Quantization/QFT & conditional/partial & MTT Born/BRST and constructive 4D
limits open \\
\bottomrule
\end{tabular}
\end{center}

\section{Strict-upgrade program}

The baseline is not reopened by stronger goals. The nine upgrades are:
\begin{enumerate}
\item[U1.] zero-primitive empirical source: open;
\item[U2.] literal global Cech--HYM/QaSU3: closed 2/2;
\item[U3.] official joint input likelihood: partial;
\item[U4.] CKM prediction beyond source rows: closed at profile criterion;
\item[U5.] absolute neutrino mass and ontology: partial;
\item[U6.] strong-CP selection: partial;
\item[U7.] MTT-derived quantization: partial;
\item[U8.] constructive nonperturbative four-dimensional QFT: partial; and
\item[U9.] unique observed-branch selection: partial at orientation level.
\end{enumerate}

\section{Next computational milestones}

\begin{enumerate}
\item \textbf{Neutral physical completion.} Select either a non-affine
spectral-action slope plus one universal scale, or dimensionful
Dirac--Majorana/seesaw blocks. Common scaling alone is proved impossible.
\item \textbf{Strong CP.} Construct the selected flux/threshold axion-current
map that changes the cancelling complete-$27$ anomaly into a justified low-
energy effective anomaly, then test quality and EDM bounds.
\item \textbf{Strict scalar source.} Emit the nine charged magnitudes,
$\lambda_H$, and physical threshold/mass rows from selected source data, or
retain them permanently as declared profile coordinates.
\item \textbf{Branch uniqueness.} Upgrade the selected retarded orientation
representative to a global action/measure and carrier uniqueness theorem.
\item \textbf{Foundational QFT.} Derive Born/record and gauge-orbit rules and
construct the required nonperturbative infinite-volume four-dimensional
limit.
\item \textbf{Zero-primitive synthesis.} Attempt U1 only after its dependency
rows close; do not replace missing source values with another calibration.
\end{enumerate}

\section{Reproducibility contract}

Every milestone must ship raw inputs, immutable source identifiers, units,
scheme/scale/loop order, branch choices, parameter ledger, covariance,
executable code, deterministic outputs, unit tests, and an acceptance
certificate. Held-out claims additionally require a data-separation
certificate. Historical packets remain provenance and cannot override their
audited successors.

\section{Conclusion}

The program has advanced substantially beyond the original roadmap, but not in
the manner that document claimed. The strongest current achievement is a
reproducible, shared-operator embedding of the renormalized Standard Model at a
clearly declared profile standard, together with exact finite/HYM and CKM
strict upgrades. The old Tier-3 and Calabi--Yau numerical chain is retired.

Progress toward a deeper theory is now sharply measurable: close U5, U6,
U7--U9 and finally U1 with source theorems rather than fitted replacements.
This roadmap is intentionally non-looping: a row changes status only when its
named missing object passes its stated acceptance test.

\begin{thebibliography}{99}

\bibitem{MTT-Admissibility}
P.~Nero,
\emph{Modal Triplet Theory: Admissibility, Encodings, and the Structure of Physical Description},
Zenodo preprint, January 2026.
\url{https://doi.org/10.5281/zenodo.18255621}

\bibitem{MTT-Foundation}
P.~Nero,
\emph{Modal Triplet Theory: Foundation},
Zenodo preprint, September 2025.
\url{https://doi.org/10.5281/zenodo.16949762}

\bibitem{MTT-FixedPoints}
P.~Nero,
\emph{Fixed Points I--VI: Complete Coherence Spine},
Zenodo preprints, August 2025.
\url{https://doi.org/10.5281/zenodo.16948748}

\bibitem{MTT-ProjectionAdmissibility}
P.~Nero,
\emph{The Projection--Admissibility Principle: Structural Constraints on Effective Physical Description},
Zenodo preprint, January 2026.
\url{https://doi.org/10.5281/zenodo.18255838}

\bibitem{MTT-Closure}
P.~Nero,
\emph{Closure and Inevitability in Modal Triplet Theory},
Zenodo preprint, January 2026.
\url{https://doi.org/10.5281/zenodo.18255510}

\bibitem{MTT-CoherenceCapacity}
P.~Nero,
\emph{Coherence Capacity as the Fundamental Resource of Effective Physics},
Zenodo preprint, January 2026.
\url{https://doi.org/10.5281/zenodo.18255905}

\bibitem{MTT-CoherenceDynamics}
P.~Nero,
\emph{Dynamics of Coherence Capacity: Transport, Concentration, and Exhaustion},
Zenodo preprint, January 2026.
\url{https://doi.org/10.5281/zenodo.18256048}

\bibitem{MTT-QM}
P.~Nero,
\emph{Modal Triplet Theory: From MTT to Quantum Mechanics},
Zenodo preprint, September 2025.
\url{https://doi.org/10.5281/zenodo.17074246}

\bibitem{MTT-QFT}
P.~Nero,
\emph{From MTT to Quantum Field Theory},
Zenodo preprint, 2025.
\url{https://doi.org/10.5281/zenodo.17068816}

\bibitem{MTT-GR}
P.~Nero,
\emph{Modal Triplet Theory: From MTT to General Relativity},
Zenodo preprint, October 2025.
\url{https://doi.org/10.5281/zenodo.16950597}

\bibitem{MTT-UVQG}
P.~Nero,
\emph{Modal Triplet Theory: From MTT to a UV-Finite, Unitary Quantum Gravity},
Zenodo preprint, 2025.
\url{https://doi.org/10.5281/zenodo.17077671}

\bibitem{MTT-Measurement}
P.~Nero,
\emph{Measurement as Disturbance and Stabilization in Modal Triplet Theory},
Zenodo preprint, 2025.
\url{https://doi.org/10.5281/zenodo.17177404}

\bibitem{MTT-Irreversibility}
P.~Nero,
\emph{Projection, Probability, and Irreversibility: Shadow Bridges Between Measurement, Black Holes, and Cosmology in Modal Triplet Theory},
Zenodo preprint, January 2026.
\url{https://doi.org/10.5281/zenodo.18256408}

\bibitem{MTT-Bell-Beables}
P.~Nero,
\emph{Modal Fixed Points, Bell’s Beables, and the Limits of Factorization},
Zenodo preprint, 2025.
\url{https://doi.org/10.5281/zenodo.17076300}

\bibitem{MTT-Bell-Temporal}
P.~Nero,
\emph{Temporal Bell Inequalities and Global Consistency in Modal Triplet Theory},
Zenodo preprint, August 2025.
\url{https://doi.org/10.5281/zenodo.18208884}

\bibitem{MTT-Stochastic}
P.~Nero,
\emph{From Modal Triplet Theory to Indivisible Stochastic Processes: A First-Principles, Fully Rigorous Derivation},
Zenodo preprint, January 2026.
\url{https://doi.org/10.5281/zenodo.18254862}

\end{thebibliography}

\end{document}
